\documentclass{amsproc} 
\usepackage{subfiles}
\usepackage[utf8]{inputenc}
\usepackage[english]{babel}
\usepackage{comment}
\usepackage[alphabetic]{amsrefs}
\usepackage{tikz}
\usepackage{xcolor}
\usepackage{datetime2} 
\usepackage[colorlinks=true,linkcolor=blue,]{hyperref}
\usepackage{tcolorbox}
\usepackage{xstring}

\usepackage{stackengine}
\usepackage{scalerel}
\usepackage{mathtools}
\usepackage{calc}
\usepackage{amsthm}
\usepackage{thmtools}
\usepackage[framemethod=TikZ]{mdframed}
\usepackage{amssymb}
\usepackage{amsfonts}
\usepackage{euscript}
% \def\EuScript#1{\mathscr#1}
% \usepackage{fourier-otf}
\usepackage{fourier}
%\usepackage{palatino}
% \usepackage[sc]{mathpazo} % add possibly `sc` and `osf` options
%\usepackage{eulervm}
%\usepackage{bbm}
%\usepackage{latexsym}
% \usepackage{mathrsfs}
% \usepackage{stmaryrd}
\def\nonfork{\mathrel{\raise0.2ex\hbox{\ooalign{\hidewidth$\vert$\hidewidth\cr\raise-0.8ex\hbox{$\smile$}}}}}

% \def\wneg{\stackon[-.2ex]{$\neg$}{$\neg$}}
\usepackage{dsfont}
% \newcommand*{\TakeFourierOrnament}[1]{{%
% \fontencoding{U}\fontfamily{futs}\selectfont\char#1}}
% \renewcommand*{\danger}{\TakeFourierOrnament{66}}
\parindent0ex
\parskip1.2ex
% \makeatletter
% \DeclareOldFontCommand{\rm}{\normalfont\rmfamily}{\mathrm}
% \DeclareOldFontCommand{\sf}{\normalfont\sffamily}{\mathsf}
% \DeclareOldFontCommand{\tt}{\normalfont\ttfamily}{\mathtt}
% \DeclareOldFontCommand{\bf}{\normalfont\bfseries}{\mathbf}
% \DeclareOldFontCommand{\it}{\normalfont\itshape}{\mathit}
% \DeclareOldFontCommand{\sl}{\normalfont\slshape}{\@nomath\sl}
% \DeclareOldFontCommand{\sc}{\normalfont\scshape}{\@nomath\sc}
% \makeatother

\def\P{\EuScript P}
\def\D{\EuScript D}
\def\A{\EuScript A}
\def\E{\EuScript E}
\def\T{\EuScript T}
\def\X{\EuScript X}
\def\Y{\EuScript Y}
\def\Z{\EuScript Z}
\def\C{\EuScript C} 
\def\U{\EuScript U}
\def\Hh{\EuScript H}
\def\I{\EuScript I}
\def\J{\EuScript J}
\def\V{\EuScript V}
\def\W{\EuScript W}
\def\R{\EuScript R}
\def\F{\EuScript F}
\def\G{\EuScript G}
\def\B{\EuScript B}
\def\M{\EuScript M}
\def\N{\EuScript N}
\def\phi{\varphi}
\def\theta{\vartheta}

\def\LDelta{L\kern-.4ex^{\scriptscriptstyle\Delta}}
\def\UDelta{\U\kern-.2ex^{\scriptscriptstyle\Delta}}
\def\BDelta{\smash{\Delta\kern-.4ex\raisebox{0.9ex}{\sf\tiny B}}\kern-.1ex}

\newcommand{\mylabel}[1]{{#1}\hfill}
\renewenvironment{itemize}
  {\begin{list}{$\triangleright$}{%
  \setlength{\parskip}{0mm}
  \setlength{\topsep}{.1\baselineskip}
  \setlength{\rightmargin}{0mm}
  \setlength{\listparindent}{0mm}
  \setlength{\itemindent}{0mm}
  \setlength{\labelwidth}{3ex}
  \setlength{\itemsep}{.1\baselineskip}
  \setlength{\parsep}{.1\baselineskip}
  \setlength{\partopsep}{0mm}
  \setlength{\labelsep}{1ex}
  \setlength{\leftmargin}{\labelwidth+\labelsep}
  \let\makelabel\mylabel}}{%
\end{list}}


\newtheoremstyle{mio}% name
     {2\parskip}     % Space above
     {2\parskip}     % Space below
     {}%         Body font
     {}%         Indent amount (empty = no indent, \parindent = para indent)
     {\bfseries}% Thm head font
     {}%        Punctuation after thm head
     {1ex}%     Space after thm head: " " = normal interword space;
           %   \newline = linebreak
     {\llap{\thmnumber{#2}\hskip0.9ex}\thmname{#1}\thmnote{\bfseries{}#3}}% Thm head spec (can be left empty, meaning `normal')

\newcounter{thm}

\tcbset{
  mythm/.style={
    colback=black!5!white,
    colframe=white!50!black,
    extrude right by=2ex,
    extrude left by=6.5ex,
    % before=\par\noindent,
    % after=\par\noindent,
    top=0.7ex,
    bottom=1.5ex,
    right=2ex,
    left=6.5ex,
    boxsep=0pt,
    % parbox=false,
    before skip=\baselineskip,
    after skip=\baselineskip,
  },
}
\theoremstyle{mio}
\newtheorem{theorem}[thm]{Theorem}\tcolorboxenvironment{theorem}{mythm}
\newtheorem{corollary}[thm]{Corollary}\tcolorboxenvironment{corollary}{mythm}
\newtheorem{proposition}[thm]{Proposition}\tcolorboxenvironment{proposition}{mythm}
\newtheorem{lemma}[thm]{Lemma}\tcolorboxenvironment{lemma}{mythm}
\newtheorem{fact}[thm]{Fact}\tcolorboxenvironment{fact}{mythm}
\newtheorem{definition}[thm]{Definition}\tcolorboxenvironment{definition}{mythm}
\newtheorem{assumption}[thm]{Assumption}\tcolorboxenvironment{assumption}{mythm}
\newtheorem{void}[thm]{}\tcolorboxenvironment{void}{mythm}
\newtheorem{notation}[thm]{Notation}\tcolorboxenvironment{notation}{mythm}
\newtheorem{note}[thm]{Note}\tcolorboxenvironment{note}{mythm}
\newtheorem{remark}[thm]{Remark}\tcolorboxenvironment{remark}{mythm}
\newtheorem{definition_theorem}[thm]{Definition\,/\,Theorem}\tcolorboxenvironment{definition_theorem}{mythm}
\newtheorem{exercise}[thm]{Exercise}
\newtheorem{example}[thm]{Example}\tcolorboxenvironment{example}{mythm}
\newtheorem{question}[thm]{Question}


\makeatletter
\providecommand{\proofNameStyle}{\bfseries}
\renewenvironment{proof}[1][\proofname]{\par
  \pushQED{\qed}%
  \normalfont% \topsep6\p@\@plus6\p@\relax
  % \vspace*{-\baselineskip}
  \trivlist
  \item[\hskip\labelsep
        \proofNameStyle
    #1\@addpunct{.}]\ignorespaces
}{%
  \popQED\endtrivlist\@endpefalse
}
\makeatother



\renewcommand*{\emph}[1]{%
   \smash{\tikz[baseline]\node[rectangle, fill=teal!25, rounded corners, inner xsep=0.5ex, inner ysep=0.2ex, anchor=base, minimum height = 2.7ex]{\strut #1};}}

%%%%%%% GETCOMMIT
\newcommand\dotGitHEAD{}
\newcommand\branch{}


\makeatletter\let\myfilehandle\@inputcheck\makeatother

\openin\myfilehandle=.git/HEAD\relax

\begingroup\endlinechar-1
  \global\read\myfilehandle to \dotGitHEAD
\endgroup
\closein\myfilehandle

\newcommand\GetBranch{}
\def\GetBranch ref: refs/heads/#1\relax{\renewcommand{\branch}{#1}}

\expandafter\GetBranch\dotGitHEAD\relax

\openin\myfilehandle=.git/refs/heads/\branch\relax

\begingroup\endlinechar-1
  \global\read\myfilehandle to \commit
\endgroup
\closein\myfilehandle

\makeatother

\linespread{1.1}
%\author{}
\author{Domenico Zambella}
% \author{al.}
\thanks{Dipartimento di Matematica, Universit\`a di Torino, via Carlo Alberto 10, 10123 Torino.}
\begin{document}
\title{After Simon, after Hrushovski}
% \hfill\texttt{~}
% \hfill\texttt{Branch \branch\ Last\_commit \StrLeft{\commit}{7}\ \DTMnow}\par
\maketitle
\raggedbottom

\begin{abstract}\parindent0ex\parskip1ex
  This is a rewriting of~\cite{Si} and~\cite{Hr}.
  But I took a few liberties in the setting and the notation.
  
  E.g.\@ I use a more compact notation when defining the patterns language.

  I work with realizations of types over $M$ instead of types themselves, and I use the notation $\N$ for a suitable set of realizations of the types over $M$.

  I restrict the formulas used to define the pattern language because I would like to accomnodate the example in Section~\ref{semigroups}~--~this only partially succeeds.

  I replace the formulas $\alpha(z)$ with types (and coheirs with strong coheirs) as presented in appendix A of~\cite{Hr}.
\end{abstract}


\section{Patterns}
\label{patterns}

\def\ceq#1#2#3{\noindent\parbox[t]{26ex}{$\displaystyle #1$}\parbox{5ex}{\hfil $#2$}{$\displaystyle #3$}}

In the following $\U$ is some monster model of a complete theory $T$ of language $L$.
For simplicity $L$ is assumed to be one-sorted.

We define a new structure $\N$ whose domain is a subset of $\U^x$, where $x$ is some fixed tuple of variables (as a matter of fact, we could safely assume $|x|=1$).
The language of $\N$, which is denoted by $L^{\rm\! p}$, is defined below.
First we need some notation.

Let $z=\langle z_i:i<\lambda\rangle$ be a tuple of (so-called) parameter variables.
We will work with a fixed set of formulas $L_z\subseteq\Delta(x\,;z)\subseteq L_{x,z}$.
Let $V$ be some large set of variables that are copies of $x$.
Finally, let $\Delta$ be the set of Boolean combinations of formulas in $\ \bigcup_{x'\in V}\Delta(x'\,;z)$.
The choice of $\Delta$ is justified in Remark~\ref{rem_first}.

Let $A\subseteq\U^x$ and $B\subseteq\U^z$.
We write $A,\Delta$ for the set of formulas $\phi(\bar a\,;z)$ for $\phi(\bar x\,;z)\in\Delta$ and $\bar a\in A^{\bar x}$, and we write $\Delta,B$ for the set of formulas $\phi(\bar x\,;b)$ for $\phi(\bar x\,;z)\in\Delta$ and $b\in B^z$.

Let $\bar a,\bar a'\in\big(\U^x\big)^n$ and  $b,b'\in\U^z$.
If $\phi(\bar a\,;b)\leftrightarrow\phi(\bar a'\,;b)$ holds for every $\phi(\bar x,z)\in\Delta$ and $b\in B^z$, we write \emph{$\bar a\equiv_{\Delta,B}\bar a'$.}
We define \emph{$b\equiv_{A,\Delta}b'$} similarly.
%We write $\Delta\mbox{-tp}(\bar a/M)$ for the set of formulas $\phi(\bar x\,;b)$ $\phi(\bar x,z)\in\Delta$ and $b\in M^z$.
% or equivalently $\Delta\mbox{-tp}(\bar a/M)=\Delta\mbox{-tp}(\bar a'/M)$

In this notes, by a \emph{finitary type\/} we mean a parameter-free type where only finitely many variables (among those displayed) occur.
For every formula $\phi(x_0,\dots,x_{n-1}\,;z)\in\Delta$ and every finitary type $\alpha(z)\subseteq L$ the language $L^{\rm\! p}$ contains the $n$-ary relation \emph{$r_{\phi\,;\alpha}$.}
We abbreviate $r_{\phi\,;\top}$ with \emph{$r_{\phi}$.}
Indeed, the reader is invited to read $\alpha(z)$ as $\top$ untill Section~\ref{Ellis}, where it plays an essential role.

Below the symbol \emph{$L^{\rm\! p}$} will also denote the set of \textit{atomic\/} formulas $r_{\phi\,;\alpha}(x_0,\dots,x_{n-1})$ for any $x_0,\dots,x_{n-1}\in V$.
Note that equalities are \textit{not\/} included in $L^{\rm\! p}$.

Given a pair of models $M\preceq N$, where $M$ is $\omega$-saturated and $\omega$-homogeneous and $N$ is $|M|^+$-saturated, we define a model $\N$ with domain $N^x$.
% Eventually, we will proved that the main definition in this section, the definition of \textit{core,} does not depend on the choice of $M$.
% However, some of the intermediate notions introduced here do depend on $M$, though to avoid cluttering the notation we will omit reference to it.
We define the interpretation of $r_{\phi\,;\alpha}$ in $\N$.
If $\bar a=a_0,\dots,a_{n-1}\in N^x$, we declare $r_{\phi\,;\alpha}(\bar a)$ true if 

\ceq{\hfill r_{\phi;\alpha}(\bar a)}{\Leftrightarrow}{\alpha(M^z)\ \cap\ \phi(\bar a\,;M^z)\ =\ \varnothing.}

We say that $\phi'(\bar x\,;z)$ is an \emph{alphabetic variant\/} of the formula $\phi(\bar x\,;z)$ if it is obtained by replacing the $z$-variables effectively occurring in $\phi(\bar x\,;z)$ by other variables among those in $z$.
Similarly, for finitary types.
Clearly, if $\phi'(\bar x\,;z)$ and $\alpha'(z)$ are alphabetic variants of $\phi(\bar x\,;z)$ and $\alpha(z)$ respectively, then $r_{\phi;\alpha}(\bar x)$ and $r_{\phi';\alpha'}(\bar x)$ are equivalent.

The following is paramount.

\begin{remark}\label{rem_first}
  By the choice of $\Delta$, the truth of $r_{\phi;\alpha}(\bar a)$ only depends on the equivalence class of $\equiv_{\Delta,M}$ of the \textit{components\/} of the tuple $\bar a$.
  Explicitely, if $\bar a=a_0,\dots,a_{n-1}$ and $\bar a'=a'_0,\dots,a'_{n-1}$ are such that $a_i\equiv_{\Delta,M} a'_i$ for every $i<n$, then \smallskip
  
  \ceq{\#\hfill r_{\phi;\alpha}(\bar a)}{\leftrightarrow}{r_{\phi;\alpha}(\bar a').} \smallskip

\end{remark}

By the remark, it would be natural to work in the quotient $\N/{\equiv_{\Delta,M}}$.
But we prefer to keep the original domain $\N$ to have a common playground with $L$.

The set of maximally consistent set of formulas of the form $\phi(x\,;b)$, where $\phi(x\,;z)\in\Delta$ and $b\in M^z$ is denoted by  $S_{x,\Delta}(M)$.
We will often confuse it with either $\U^x/{\equiv_{\Delta,M}}$ or $\N/{\equiv_{\Delta,M}}$, with a preference for the latter.

A \emph{p-type\/} is a set of formulas in $L^{\rm\! p}(\N)$. 

% We write \emph{$\textrm{p-tp}(\bar a)$} for the p-type of $\bar a$, i.e.\@ the set of formulas $r_{\phi}(\bar x)$ such that $r_{\phi}(\bar a)$ holds.

Note incidentally that the relation $x_0\equiv_{\Delta,M} x_1$ is definable by a p-type.
Namely the p-type containing the formulas $r_{\theta}(x_0,x_1)$ for every $\theta(x_0,x_1\,;z)$ of the form $\phi(x_0\,;z)\nleftrightarrow\phi(x_1\,;z)$.

\begin{fact}\label{fact_p_qe}
  Formulas in $L^{\rm\! p}$ are closed under adjunction of dummy variables, disjunction, conjunction, and existential quantification.
\end{fact}

\begin{proof}
  The claim about dummy variables is clear, so we tacitatly use it below.

  Up to replacing $\phi(\bar x\,;z)$ and $\alpha(z)$ with alphabetic variants, we can assume that the $z$-variables that effectively occur in $\phi(\bar x\,;z)$ and $\alpha(z)$ are disjoint from those that occur in $\psi(\bar x\,;z)$ and $\beta(z)$ .
  Therefore the following equivalences hold 

  \ceq{\hfill r_{\phi\,;\alpha}(\bar a)\vee r_{\psi\,;\beta}(\bar a)}{\Leftrightarrow}{\forall z\in\alpha(M^z)\neg \phi(\bar a\,;z)\ \ \vee\ \ \forall z\in\beta(M^z)\neg \psi(\bar a\,;z)} 
  
  \ceq{}{\Leftrightarrow}{\forall z\in\alpha(M^z)\cap\beta(M^z)\ \ \neg \phi(\bar a\,;z)\, \vee\,\neg\psi(\bar a\,;z)} 

  % \ceq{}{\Leftrightarrow}{r_{\phi\vee\psi\,;\,\alpha\cap\beta}(\bar a)}
  
  % \ceq{\hfill r_{\phi}(\bar a)\vee r_{\psi}(\bar a)}{\Leftrightarrow}{\forall z\in M^z\ \phi(\bar a\,;z)\ \ \vee\ \ \forall z\in M^z\ \psi(\bar a\,;z)}

  % \ceq{}{\Leftrightarrow}{\forall z\in M^z\ \big[ \ \phi(\bar a\,;z)\, \vee\,\psi(\bar a\,;z)\big]}  

  % \ceq{}{\Leftrightarrow}{r_{\phi\vee\psi}(\bar a)}

  A similar argument applies to the conjunction of two formulas.

  Finally, consider the formula $\exists y\,r_{\phi\,;\alpha}(\bar x,y)$.
  By the definition of $\Delta$, we can assume that $\phi$ has the form

  \ceq{\hfill\phi(\bar x,y\,;z)}{=}{\bigvee_{i=1}^m\phi_i(\bar x,y\,;z)}

  where $\phi_i(\bar x,y\,;z)\,=\,\phi'_i(\bar x\,;z)\,\wedge\,\psi_i(y\,;z)$ and each $\phi'_i(\bar x\,;z)$ is a conjunction of formulas in $\bigcup_{x'\in V} L_{x'\,;z}$ or negation thereof.
  Closure under existential quantification follows from the following equivalences

  \ceq{\hfill\exists y\ r_{\phi\,;\alpha}(\bar a, y)}{\Leftrightarrow}{\exists y\,\forall z\Big[ z\in\alpha(M^z)\rightarrow \neg\bigvee_{i=1}^m\phi_i(\bar a,y\,;z)\Big]}

  \ceq{}{\Leftrightarrow}{\forall z\Big[ z\in\alpha(M^z)\rightarrow \neg\Big(\bigvee_{i=1}^m\phi'_i(\bar a\,;z)\wedge\forall y\,\psi_i(y\,;z)\Big)\Big]}

  \ceq{}{\Leftrightarrow}{\vphantom{\bigvee_{i=1}^n}r_{\theta\,;\alpha}(\bar a),}

  % \ceq{\hfill\exists y\ r_{\phi}(\bar a, y)}{\Leftrightarrow}{\exists y\,\forall z\in M^z \bigvee_{i=1}^m\phi_i(\bar a,y\,;z)}

  % \ceq{}{\Leftrightarrow}{\forall z\in M^z \bigvee_{i=1}^m\phi'_i(\bar a\,;z)\wedge\exists y\,\psi_i(y\,;z)\Big]}

  % \ceq{}{\Leftrightarrow}{\vphantom{\bigvee_{i=1}^n}r_{\theta}(\bar a),}

  where 
  
  \ceq{\hfill\theta(\bar x\,;z)}{=}{\bigvee_{i=1}^m\phi'_i(\bar x\,;z)\ \wedge\ \forall y\,\psi_i(y\,;z).}

  As $\exists y\,\psi_i(y\,;z)\in L_z\subseteq\Delta$, the above formula is in $\Delta$ as required.
\end{proof}

\begin{fact}[ ($\infty$-p-saturation)]\label{fact_psat}
  Let $A\subseteq \N$ have arbitrary cardinality.
  Let $\bar x$ be a tuple of copies of $x$.
  Then every p-type $p(\bar x)\subseteq L^{\rm\! p}(A)$ that is finitely consistent in $\N$ is realized in $\N$.
\end{fact}

\begin{proof}
  Let $A'\subseteq\N$ be a small set containing a representative of every $\equiv_{\Delta,M}$ equivalence class that intersects $A$.
  Replacing the elements of $A$ that occur in $p(\bar x)$ by elements of $A'$ we obtain an equivalent type $p'(\bar x)\subseteq L^{\rm\! p}(A')$.
  Note that every formula $r_{\phi\,;\alpha}(\bar x)$ there is a $\Delta$-type $q(\bar x)\subseteq L(M)$ such that

  \ceq{\hfill \N\models r_{\phi\,;\alpha}(\bar a)}{\Leftrightarrow}{N\models q(\bar a)}

  Therefore the saturation of $N$ implies that $p'(\bar x)$ is realized in $\N$.
\end{proof}

%%%%%%%%%%%%%%%%%%%%
%%%%%%%%%%%%%%%%%%%%
%%%%%%%%%%%%%%%%%%%%
%%%%%%%%%%%%%%%%%%%%
%%%%%%%%%%%%%%%%%%%%
%%%%%%%%%%%%%%%%%%%%
\section{Strong coheirs}

We need a notion of independence that is stronger than the usual notion of coheir.
The notion is adapted from that is \cite{KNP}.
We use parameter-free types (finitary types are assumed to be parameter-free) where \cite{KNP} uses finitely many parameters from $M$.
We beleave the difference is minor but it simplifies the exposition.

Let $p(z)\subseteq L(N)$.
We say that $p(z)\subseteq L(N)$ is \emph{strongly finitely satisfiable\/} in $M$ if, for every finitary type $\alpha(z)\subseteq p$, and every formula $\phi(z)\in p$, the type $\alpha(z)\wedge\phi(z)$ is satisfied in $M$.
Recall that a finitary type is one which is partameter-free and with finitely many variables.

Let $b\in\U^z$ and $C\subseteq N$.
When ${\rm tp}(b/C)$ is strongly finitely satisfiable in $M$, we write \emph{$b\nonfork_M C$.}
Explicitely, $b\nonfork_M C$ holds when $\alpha(M^z)\cap\phi(M^z)\neq\varnothing$ for every $\phi(z)\in L(C)$ and every finitary $\alpha(z)$ such that $b\models\alpha(z)\wedge\phi(z)$.

\begin{fact}
  % Assume $M$ is $\omega$-saturated.
  Every type $p(z)\subseteq L(N)$ that is strongly finitely satifiable in $M$ has an extension to a complete type $q(z)\in S(N)$ that is strongly finitely satisfiable in $M$.
\end{fact}

\begin{proof}
  Let $q(z)\subseteq L(N)$ be maximal among the types that contain $p(z)$ and are strongly finitely satisfiable in $M$.
  We prove that $q(z)$ is complete.
  Suppose not, say it does not contain $\phi(z)$ nor $\neg\phi(z)$.
  Then there is a finitary type $\alpha(z)\subseteq p$ and a formula $\psi(z)\in p$ such that neither $\alpha(z)\wedge\psi(z)\wedge\phi(z)$ nor $\alpha(z)\wedge\psi(z)\wedge\neg\phi(z)$.
  This contradicts the strong finite satisfiability of $q(z)$.
\end{proof}

Let $\hat m$ be an enumeration of $M$ and assume that $|z|=|\hat m|$.
We define

\ceq{\hfill\emph{$\B$}}{=}{\Big\{b\in\U^z\ :\ \hat m\equiv b\nonfork_MN\Big\},}

where $\equiv$ is the elementary equivalence in $L$ (however, as $L_z\subseteq\Delta$, we could replace it with $\equiv_{\varnothing,\Delta}$).
The $N,\Delta$-topology on $\B$ is the one generated by the closed sets $\{b\in\B\ :\ \phi(\bar a\,;b)\}$ for $\phi(\bar x\,;z)\in\Delta$ and $\bar a\in N^{\bar x}$.
The $N,\Delta$-topology on $\B/{\equiv_{N,\Delta}}$ is the quotient of the topology on $\B$.

%%%%%%%%%%%%%%%%%%%%
%%%%%%%%%%%%%%%%%%%%
%%%%%%%%%%%%%%%%%%%%
%%%%%%%%%%%%%%%%%%%%
%%%%%%%%%%%%%%%%%%%%
%%%%%%%%%%%%%%%%%%%%
\section{Morphism}

A p-morphism is a map $f:\N\to\N$ such that for every $r_{\phi\,;\alpha}(\bar x)\in L^{\rm p}$ and every $\bar a\in({\rm dom}f)^{\bar x}$

\ceq{\#\hfill\N\models r_{\phi\,;\alpha}(\bar a)}{\Rightarrow}{\N\models r_{\phi\,;\alpha}(f\bar a)}.

As $\equiv_{\Delta,M}$ is definable by a p-type, p-morphism induce well-defined maps on the quotient $\N/{\equiv_{\Delta,M}}$.
We say that two p-morphism $f$ and $g$ are \emph{equivalent,} and write $f\sim g$, if they induce the same map on $\N/{\equiv_{\Delta,M}}$.
We often confuse equivalent p-morphism and we use the same symbol $f$ to denote the map induced on the quotient.

An \emph{endomorphism\/} is a total morphism.
We write \emph{${\rm End}(\N)$} for the set of p-endo\-mor\-ph\-isms of $\N$.

\begin{definition}\label{def_topology}
  We endow ${\rm End}(\N)$ with the topology that has as basic closed sets those of the form 

  \ceq{\hfill\Big\{\ f\in{\rm End}(\N)}{:}{\N\ \models\ r_{\phi\,;\alpha}\big(f\!a_1,\dots,f\!a_n,c_1,\dots,c_m\big)\Big\}}\smallskip

  for $\phi(x\,;z)\in\Delta$ and $a_1,\dots,a_n,c_1,\dots,c_m\in N^x$.
\end{definition}

Note that, by Remark~\ref{rem_first}, if $f\sim g$ then $f$ and $g$ belong to the same basic closed sets, therefore this topologies naturally induces a topology on ${\rm End}(\N)/{\sim}$.

\begin{fact}
  With the above topology ${\rm End}(\N)$ is T$_1$ (i.e.\@ singletons are closed) and compact.
  Composition of p-morphism, which we denote by $*$, is a left continuous operation.
\end{fact}

\begin{proof}
  We prove that singletons are closed.
  In fact

  \ceq{\hfill\big\{f\big\}}{=}{\bigcap_{a\in N^x}\Big\{ g\in{\rm End}(\N)\ :\ a\equiv_{\Delta,M}ga\Big\}}

  therefore, as $\equiv_{\Delta,M}$ is defined by a p-type, $\{f\}$ is the intersection of closed sets.
  
  We prove compactness.
  Assume the following intersection is finitely consistent

  \ceq{\#.\hfill}{~}{\bigcap_{i<\lambda}\Big\{\ f\in{\rm End}(\N)\ :\ \N\ \models\ r_{\phi_i;\alpha_i}\big(f\!\bar a,\bar c\big)\Big\}.}

  then the p-type

  \ceq{\hfill\bar x\equiv_{\Delta, M}\bar a}{\cup}{\big\{r_{\phi_i;\alpha_i}\big(\bar x,\bar c\big) \ :\ i<\lambda\big\}}

  is finitely consistent in $\N$.
  By $\infty$-p-saturation, it is realized by some $\bar a'$.
  By Fact~\ref{fact_phogeneity} there is a p-endomorphism $\N$ that maps $\bar a$ to $\bar a'$.
  This proves that the set in $\#$ is non empty.

  We prove left continuity of composition.
  That is for any given $f\in{\rm End}(\N)$ we prove that the function $\tau_f:{\rm End}(\N)\to{\rm End}(\N)$ that maps $g\mapsto g*f$ is continuous.
  Pick any basic closed set $\{h\, :\, r_{\phi\,\alpha}(h\bar a,\bar c)\}$.
  Then 
  
  \ceq{\hfill\tau_f^{-1}\big[\big\{h\,:\,r_{\phi\,;\alpha}\big(h\bar a,\bar c\big)\big\}\big]}{=}{\big\{g\ :\ r_{\phi\,;\alpha}\big(g(f\!\bar a),\bar c\big)\big\}}
  
  is closed, as required.
\end{proof}

If the converse of the implication in \# also holds, $f$ is called a \emph{partial p-isomorphism.}
For the back-and-forth argument in the following proof we need to take inverse of partial p-isomorphisms.
Strictly speaking partial p-isomorphism need not be invertible but for our puposes it suffices to have a p-morphism $g$ such that $f\circ g$ and $g\circ f$ are equivalent to the identity on ${\rm range}f$, respectively ${\rm dom}f$. 

% Though the switching between the two perspectives is intuitively obvious, it requires a fair amount of notation.

% Let $R\mathrel{\subseteq}\U^x\times\U^x$ be a binary relation.
% If $\bar a=\langle a_i:i<\lambda\rangle$ and $\bar b=\langle b_i:i<\lambda\rangle$ we write $\langle \bar a,\bar b\rangle\in R$ to abbreviate: $\langle a_i,b_i\rangle\in R$ for all $i<\lambda$.
% We say that $R$ is a \emph{p-morphism} if the following condition holds 

% \ceq{\#\hfill\U^x\models r_{\phi}(\bar a)}{\Rightarrow}{\U^x\models r_{\phi}(\bar b)}\hfill for every $r_{\phi}\in L^{\rm p}$ and every $\langle \bar a,\bar b\rangle\in R$.

% When the converse implication in \# also holds, i.e.\@ when also $R^{-1}$ is a p-morphism, we say that $R$ is a \emph{partial p-isomorphism.}
% When $\A\subseteq{\rm dom}\,R$, $\B\subseteq{\rm range}\,R$, and $R\cap (\A\times\B)$ is a partial p-isomorphism, we say that $R$ is an \emph{p-isomorphism} between $\A$ and $\B$ or, when $\A=\B$, a \emph{p-automorphism} of $\A$.

\begin{fact}\label{fact_phogeneity}
  Every p-morphism extends to a p-morphism defined on the whole of $\N$.
  Also, every p-isomorphism extends to a p-automorphism of $\N$.
\end{fact}

\begin{proof}
  Let $f:\N\to\N$ be a p-morphism.
  Let $\bar a$ be an enumeration of ${\rm dom} f$ and let $b\in N^x$.
  Let $p(\bar x)=\textrm{p-tp}(\bar a)$ and let $q(\bar x,y)=\textrm{p-tp}(\bar a,b)$.
  By $\infty$-p-saturation $\exists y\,q(\bar x,y)$ is equivalent to the type containing the existential closure of the formulas in $q(\bar x,y)$.
  By Fact~\ref{fact_p_qe}, the latter is equivalent to a p-type, which is necessarily coicides with $p(\bar x)$.
  Then $q(f\bar a,y)$ is consistent hence by $\infty$-p-saturation it is realized by some $c\in N^x$.
  Then $f\cup\{\langle b,c\rangle\}$ is a p-morphism defined in $b$ that extends $f$.

  The second claim is obtained by back-and-forth.
\end{proof}

%%%%%%%%%%%%%%%%%
%%%%%%%%%%%%%%%%%
%%%%%%%%%%%%%%%%%
%%%%%%%%%%%%%%%%%
%%%%%%%%%%%%%%%%%
\section{Duality}

Recall that $\hat m$ is an enumeration of $M$ and $|z|=|\hat m|$.
For convenient notation, $x$ is assumed to be of infinite length~--~then, up to equivalence, all formulas in $L_z(\U)$ have the form $\phi(a\,;z)$ for some $\phi(x\,;z)\in L$ and $a\in\U^x$.

\begin{remark}
  We stress a simple fact which is used in the proof below. 
  Recall that $\phi'(a\,;z)$ and $\alpha'(z)$ are alphabetic variants of $\phi(a\,;z)$ and $\alpha(z)$ if they are obtained by replacing the variables effectively occurring in $\phi(a\,;z)$ and $\alpha(z)$ by other variables among those in $z$.
  The following are clearly equivalent
  \begin{itemize}
    \item[1.] $\alpha(M^z)\cap\phi(a\,;M^z)=\varnothing$
    \item[2.] $\hat m\not\models\alpha'(z)\wedge\phi'(a\,;z)$ for every alphabetic variants of $\phi(a\,;z)$ and $\alpha(z)$ .
  \end{itemize}
\end{remark} 

\begin{fact}
  There is a continuous bijection $\B/{\equiv_{N,\Delta}}\to{\rm End}(\N)/{\sim}$.
\end{fact}

\begin{proof}
  In 1 below, we define a map $b\mapsto f_b$ from $\B$ into ${\rm End}(\N)$. In 2, we define a map $f\mapsto b_f$ from ${\rm End}(\N)$ into $\B$.
  The reader may check that these maps induce maps between $\B/{\equiv_{N,\Delta}}$ and ${\rm End}(\N)/{\sim}$ that are each other's inverse.
  In 3 we prove the continuity of the first map.

  1.\ 
  Let $b\in\B$ be given.
  We define
  
  \ceq{\hfill p(x,y,z)}{=}{\big\{\phi(y\,;z)\rightarrow\phi(x\,;\hat m)\ :\ \phi(x\,;z)\in\Delta\big\}.}

  We claim that if $a\equiv_{\Delta,M}a'$ then $p(x,a,b)=p(x,a,b)$.
  In fact, $\phi(a\,;b)\nleftrightarrow\phi(a'\,;b)$ would contradict $b\nonfork_MN$.
  Also, note that $p(x,a,b)$ is consistent, otherwise $\neg\exists x\bigwedge_i\phi_i(x\,;\hat m)$, for some finitely many formulas $\phi_i(x\,;z)\in\Delta$ such that $\phi_i(a\,;b)$.
  But this would contradict $b\equiv\hat m$.

  Therefore it is possible to define $f_b(a)$ to be some/any realization of $p(x,a,b)$.
  We prove that $f_b$ is a p-morphism.

  Let $\psi(x_1,\dots,x_n\,;z)\in\Delta$, $\alpha(z)\subseteq L$, and $a_1,\dots,a_n\in\U^x$ be such that $\neg r_{\psi\,;\alpha}(f_ba_1,\dots,f_ba_n)$.
  Then $\alpha(M^z)\cap\psi(f_ba_1,\dots,f_ba_n\,;M^z)\neq\varnothing$.
  Pick some alphabetic variants such $\alpha'(\hat m)$ and $\psi'(f_ba_1,\dots,f_ba_n\,;\hat m)$.
  From the definition of $f_b$ we obtain $\psi'(a_1,\dots,a_n\,;b)$ and, as $b\equiv\hat m$, we also have $\alpha'(b)$.
  Finally, as $b\nonfork_MN$, we conclude that $\alpha'(M^z)\cap\psi'(a_1,\dots,a_n\,;M^z)\neq\varnothing$, i.e.\@ that $\neg r_{\psi\,;\alpha}(a_1,\dots,a_n)$.
  
  
  

  % We can assume that $\hat m$ realizes some alphabetic variant of $\alpha(z)$, otherwise $r_{\psi\,;\alpha}(f_ba_1,\dots,f_ba_n)$ trivially holds. 
  % As $r_{\psi\,;\alpha}(x_1,\dots,x_n)$ does not depend on the alphabetic variants we assume $\hat m\models\alpha(z)$.
  % As $\hat m\equiv b\nonfork_MN$, we obtain $\neg\psi(a_1,\dots,a_n\,;b)$.

  % % We can assume $\alpha(M^z)\neq\varnothing$, otherwise $r_{\psi}(f\!a_1,\dots,f\!a_n)$ trivially holds.
  % Assume

  % \ceq{\hfill\psi(x_1,\dots,x_n\,;z)}{=}{\bigwedge_{k=1}^m\bigvee_{h=1}^m\phi_{h,k}(x_{i_{h,k}}\,;z)}

  % for some $i_{h,k}\in\{1,\dots,n\}$ and $\phi_{h,k}(x_{i_{h,k}}\,;z)\in\Delta$.
  % After re-indexing, from $\neg\psi(a_1,\dots,a_n\,;b)$ we obtain 
  
  % \ceq{}{~}{\bigwedge_{k=1}^m\neg\phi_k(a_k\,;b)}
  

  % \ceq{}{~}{\bigwedge_{k=1}^m\neg\phi_k(f_ba_k\,;\hat m)}


  % Therefore  $\phi_k(f_ba_k\,;\hat m)$ and all its alphabetic variants.
  % Therefore $\phi_k(f_ba_k;M^z)=M^z$, and $r_{\psi}(f_ba_1,\dots,f_ba_n)$ follows.

  2.\ 
  Let $f: \N\to \N$ be an p-epimorphism.
  Let $p(z)$ be the type containing the formulas $\phi(a\,;z)$ for $\phi(x\,;z)\in\Delta$ and $a\in N^x$ such that $\phi(f\!a\,;\hat m)$.
  Clearly $p(z)\vdash\hat m\equiv z$ and is complete over $N,\Delta$.
  We prove that $p(z)$ is strongly finitely satisfied in $M$, then $b_f$ is some/any realization of $p(z)$.

  Suppose for a contradiction $p(z)$ is not strongly finitely satisfied in $M$. Then $\alpha(M^z)\cap\psi(a_1,\dots,a_n\,;M^z)=\varnothing$, where $\alpha(z)\subseteq p$ is a finitary type and 
  
  \ceq{\hfill\psi(x_1,\dots,x_n\,;z)}{=}{\bigwedge_{i=1}^n\phi_i(x_i\,;z)}
  
  for some $\phi_i(x\,;z)\in L$ and $a_i\in N^x$ such that $\phi_i(f\!a_i\,;\hat m)$ for every $i$.
  Then $r_{\psi\,;\alpha}(a_1,\dots,a_n)$ and $\neg r_{\psi\,;\alpha}(f\!a_1,\dots,f\!a_n)$.
  This is not possible because $f$ is a p-morphism.

  % 3.\ We prove that the map defined in 1 is continuous.
  % Pick any basic closed set in $\B$, say $\big\{b\,:\,\phi(\bar a\,;b)\big\}$.
  % We claim that the set $\big\{f\,:\,\phi(\bar a\,;b_f)\big\}$ is closed in ${\rm End}(\N)$. Indeed, by the definition of $b_f$ in 1

  % \ceq{\hfill\big\{f\ :\ \phi\big(\bar a\,;b_f\big)\}}{=}{\Big\{f\ :\ \phi\big(f\bar a\,;\hat m\big)\Big\}}

  % \ceq{\#}{=}{\Big\{f\ :\ r_{\psi}\big(f\bar a,c\big)\Big\}}

  % where $\psi(\bar x,y;z)$ and $c$ are choosen so that $r_{\psi}(\bar x,c)$ is equivalent to $\phi(\bar x\,;\hat m)$.
  % As the set in $\#$ is closed in ${\rm End}(\N)$, continuity of the map $f\mapsto b_f$ follows.

  % We show how to choose $\psi(\bar x,y;z)$ and $c$.
  % Let $y$ be a copy of $x$.
  % The formula $\psi(\bar x,y;z)$ is obtained from $\phi(\bar x\,;z)$ by replacing the variables in the tuple $z$ effectively occurring in $\phi(\bar x\,;z)$ with variables in the tuple $y$.
  % Note that $z$ is dummy in $\psi(\bar x,y\,;z)$.
  % Finally, we choose $c\in M^{y}$ such that $\psi(\bar x,c\,;z)$ is equivalent to $\phi(\bar x\,;\hat m)$.

  3.\ 
  We prove that the map defined in 1 is continuous.
  Pick any basic closed set in ${\rm End}(\N)$, say $\big\{ h\, :\, r(h\bar a,\bar c)\big\}$.
  We claim that the set $\big\{b\, :\, r_{\psi\,;\alpha}(f_b\bar a,\bar c\big)\}$ is closed. Indeed, by the definition of $f_b$ in 1,

  \ceq{\hfill\big\{b\ :\ r_{\psi\,;\alpha}\big(f_b\bar a,\bar c\big)\big\}}{=}{\Big\{b\ :\ b\models \exists\bar x\ \Big[\bigwedge_{i=1}^np(x_i,a_i,z)\wedge r_{\psi\,;\alpha}\big(\bar x,\bar c\big)\Big]\Big\}}

  is closed in $\B$ because $r_{\psi\,;\alpha}\big(\bar x,\bar c\big)$ is definable by a $N,\Delta$-type.
\end{proof}

% We define a semigroup operation on $\B$.
% Namely, if $f_b,f_c:\to$ are the p-morphisms the correspond to $b,c\in\B$ we write $b{*}c$ for the element of $B$ that corresponds to the p-morphism $f_b\circ f_c$.

% On $\B/{\equiv_{\Delta,M}}$ the operation $*$ is well-defined and turns it into a semigroup.

% Let $p(z)$ and $q(z)$ be the types that correspond to the p-morphisms $f_p$, respectively $f_q$.
% Pick any sufficiently saturated model $N\succeq M$ and $c\models q(z)\restriction N$.
% Then the type $(p\otimes q)(z,c)$ corresponds to the p-morphisms $f_p\circ f_q$.

% 

% $(p\otimes q)(z,c)=\{\phi(z,c)\ \ :\ \ \phi(z,c')\in p \textrm{ for some/any } c'\models q(z)\restriction N\}$

% Then $\phi(a',c)$ holds for every $a'\in f_q^{-1}[N]$.

% $z'$ be a copy of $z$
% Say $f,g: S(M)\to S(M)$ and $h=gf$.
% Let $m,n$ enumerate $M$.
% Write $p(z)=\Delta\mbox{-tp}(m/\varnothing)$.
% By Lemma 3.14, we have that $g$ and $f$ correspond to global extensions $p_f$ and $p_g$. 
% Explicitely:

% $p_f=\{\phi(a,z)\ :\ φ(fa,m),\ a\in\U^x\}$

% $p_g=\{\phi(a,z)\ :\ φ(ga,m),\ a\in\U^x\} = \{\phi(fa,z)\ :\ φ(gfa,m),\ a\in\U^x\}$

% $p_{gf}=\{\phi(a,z)\ :\ φ(gfa,m),\ a\in\U^x\}$

% % f(tp(e/A))={φ(x,b) : φ(e',y)∈f*p for some/any e' ≡_A e}={φ(x,b) : φ(e,y)∈f*p}
% % g*q={ φ(e,z) : φ(x,c)∈ g(tp(e/B)) }
% % g(tp(e/B))={φ(x,c) : φ(e',z)∈g*q for some/any e' ≡_B e}={φ(x,c) : φ(e,z)∈g*q}

% Now, h corresponds to a global extension h*q of q. Let's write h*q in terms of f*p and g*q:
% h*q={φ(e,z) : φ(e,z)\in gf(tp(e/A))}={φ(e,z) : φ(e',z)∈g*q for some/any e' realising f(tp(e/A))}={φ(e,z) : φ(e',z)∈g*q for some/any e' such that ⊧θ(e',b) for any θ(e,y)∈f*p}={φ(e,z) : φ(e',z)∈g*q for some/any e' such that e'b≡eb' for any b' realising f*p|e}

% Hence, for p(y) global A-invariant type extending tp(b/ø) and q(z) global B-invariant type extending tp(c/ø), define
% q⊗p={φ(e,z) : φ(e',z)∈q for some/any e' such that e'b≡eb' for any b' realising p|e}.
% We can further check that q⊗p is finitely satisfiable in A if q is finitely satisfiable in B and p is finitely satisfiable in A.  


%%%%%%%%%%%%%%%%%%%%
%%%%%%%%%%%%%%%%%%%%
%%%%%%%%%%%%%%%%%%%%
%%%%%%%%%%%%%%%%%%%%
%%%%%%%%%%%%%%%%%%%%
\section{Semigroups}\label{semigroups}

Assume $\U$ is a unital semigroup (a monoid) w.r.t.\@ some definable operation which we denote by $\cdot$.
The unit is denoted by $1$.
In this section we could work with $|x|=|z|=1$ but, for uniformity with the previous section we take $|x|=|z|=\omega$.
We write $z\cdot x$ for coordinatewise multiplication.
In this section $\Delta(x\,;z)$ contains the formulas $\phi(z\cdot x)$ for all $\phi(x)\in L$.

\begin{fact}
  Assume that $a\cdot M=M$ for all $a\in M$, which holds in particular when $\U$ is a group.
  For $b\in\B$, let $g_b:\N\to\N$ be the function that maps $a\mapsto b\cdot a$.
  Then $b\mapsto g_b$ induces a bijection between $\B/{\equiv_{N,\Delta}}$ and ${\rm End}(\N)/{\sim}$.
\end{fact}

\begin{proof}
  In 1 below we prove that $g_b$ is a p-morphism.
  In 2 we define a map $g\mapsto b_g$ from ${\rm End}(\N)$ into $\B$.
  The reader may check that these maps induce maps between $\B/{\equiv_{N,\Delta}}$ and ${\rm End}(\N)/{\sim}$ that are each other's inverse.
  In 3 we prove the continuity of the map $b\mapsto g_b$.

  1.\ 
  Assume $r_{\psi\,;\alpha}(a_1,\dots,a_n)$.
  We prove $r_{\psi\,;\alpha}(g_ba_1,\dots,g_ba_n)$.
  Suppose not, then $\alpha(M^z)\cap\psi(M^z\cdot b\cdot a_1,\dots,;M^z\cdot b\cdot a_n)\neq\varnothing$.
  From  $b\nonfork_MN$, we obtain $\alpha(M^z)\cap\psi(M^z\cdot a_1,\dots,;M^z\cdot a_n)\neq\varnothing$.
  This contradicts $r_{\psi\,;\alpha}(a_1,\dots,a_n)$.
  
  2.\ 
  Let $g:\N\to\N$ be a p-morphism.
  Let $p(y)$ be the type containing $\phi(\hat m\cdot y\cdot a)$ for all $\phi(x)\in L$ and $a\in N^x$ such that $\phi(\hat m\cdot ga)$.
  We prove that $p(y)$ is strongly finitely satisfied in $M$.
  If for a contradiction $\alpha(M^z)\cap\psi(a_1,\dots,a_n\,;\hat m\cdot M^z)=\varnothing$, where 

  \ceq{\hfill\psi(x_1,\dots,x_n\,;z)}{=}{\bigwedge_{i=1}^n\phi_i(z\cdot x_i)}

  for some $\phi_i(x)\in p$ and $a_i\in N^x$ such that $\phi_i(\hat m\cdot g\!a_i)$ for every $i$.
  As $\hat m\cdot M^z=M^z$, we have that $r_\psi(a_1,\dots,a_n)$ and $\neg r_\psi(g\!a_1,\dots,g\!a_n)$.
  This is not possible because $g$ is a p-morphism.
  Then any $b_g$ that realizes $p(z)$ is such that $g\sim g_{b_g}$.
 
  3.\ 
  We prove that the map defined in 1 is continuous.
  Pick any basic closed set in ${\rm End}(\N)$, say $\big\{ h\, :\, r(h\bar a,\bar c)\big\}$ where $\bar a=a_1,\dots,a_n$ and $\bar c=c_1,\dots,c_m$.
  We claim that the set $\big\{b\, :\, r(g_b\bar a,\bar c\big)\}$ is closed.
  Indeed, by the definition of $g_b$ in 1,

  \ceq{\hfill\big\{b\ :\ r\big(g_b\bar a,\bar c\big)\big\}}{=}{\big\{b\ :\ r\big(b\cdot a_1,\dots,b\cdot a_n,\bar c\big)\big\}}
  
  is closed in $\B$ because $r\big(w\cdot a_1,\dots,w\cdot a_n,\bar c\big)$ is definable by a $N,\Delta$-type $p(w,\bar a,\bar c)$.
\end{proof}

%%%%%%%%%%%%%%%%%%%%
%%%%%%%%%%%%%%%%%%%%
%%%%%%%%%%%%%%%%%%%%
%%%%%%%%%%%%%%%%%%%%
%%%%%%%%%%%%%%%%%%%%
\section{Ellis's semigroup}\label{Ellis}

% As $N\preceq\U$ is choosen to be $|M|^+$-saturated, it is possible to identify $\N/{\equiv_{\Delta,M}}$ and $\U^x/{\equiv_M}$.
We write ${\rm Aut}(N/\{M\})$ for the group of automorphisms of $N$ that fix $M$ setwise.
In this section we show that ${\rm End}(\N)$, endowed with a topology that is finer than that introduced above, is homeomorphic to the Ellis semigroup associated to the action % of ${\rm Aut}(M)$ on $S_{x,\Delta}(M)$.
% We prefer to work with the action of 
${\rm Aut}(N/\{M\})$ on $N^x$.
Note that this is essentially equivalent to the action of ${\rm Aut}(M)$ on $S_{x,\Delta}(M)$~--~we chose the former based on personal preference.

We recall the definition of this semigroup, which we denote by ${\rm El}_x(N,M)$.
On $\N$ we assume the topology generated by the closed sets of the form $\phi(N^x;b)$ for $\phi(x\,;z)\in\Delta$ and $b\in M^z$\!.
This we call the \emph{logic topology.}
We identify ${\rm Aut}(N/\{M\})$ with a subset of $(N^x)^{N^x}$\!\!, the set of all maps from $N^x$ into itself, which we endow with the product topology.
The Ellis semigroup is the closure of ${\rm Aut}(N/\{M\})$ in the product topology.
The following theorem proves that ${\rm El}_x(N,M)$ concides with ${\rm End}(\N)$.

\begin{theorem}
  Assume $M$ is $\lambda$-homogeneous (where $\lambda=|z|$).
  Then ${\rm End}(\N)$ is a closed subset of $(N^x)^{N^x}$\!\!, and ${\rm Aut}(N/\{M\})$ is dense subset of it.
\end{theorem}

Note that the topology that ${\rm End}(\N)$ inherits from $(N^x)^{N^x}$\! is finer than the topology introduced in Definition~\ref{def_topology}.

\begin{proof}
  It is clear that ${\rm End}(\N)$ is a closed set containing ${\rm Aut}(N/\{M\})$.
  Let $f\in{\rm End}(\N)$ and $a_1,\dots,a_n\in N^x$ be given.
  Let $\phi_i(x\,;z)\in\Delta$ and $b\in M^z$ be such that $fa_i\models\phi_i(x\,;b)$.
  We need to prove that there is $g\in{\rm Aut}(N/\{M\})$ such that $ga_i\models\phi_i(x\,;b)$.
  Let $\alpha(z)={\rm tp}(b)$.
  Let

  \ceq{\hfill\psi(x_1,\dots,x_n\,;z)}{=}{\neg\bigwedge_{i=1}^n\phi_i(x_i\,;z)}

  As $\neg r_{\psi\,;\alpha}(fa_1,\dots,fa_n)$ and $f$ is a p-morphism, $\neg r_{\psi\,;\alpha}(a_1,\dots,a_n)$.
  Then there is $c\models\alpha(z)$ such that $\phi_i(a_i\,;c)$ for every $i$.
  As $M$ is $\lambda$-homogeneous there is a $g\in{\rm Aut}(N/\{M\})$ such that $gc=b$, then $\phi_i(ga_i\,;b)$ as requiured.
\end{proof}

% Let $\gamma:\U^x\to\U^x$ be a function that is approximated by automorphisms in ${\rm Aut}(\U/M)$.
% That is, for every finite set $A\subseteq\U^x$ there is an automorphism $f_A\in{\rm Aut}(\U/M)$ that $f_Aa=\gamma a$ for every $a\in A$.
% We associate to $\gamma$ the endomorphism $f_\gamma\in{\rm End}(\N)$ defined as follows.

%%%%%%%%%%%%%%%%%%%%
%%%%%%%%%%%%%%%%%%%%
%%%%%%%%%%%%%%%%%%%%
%%%%%%%%%%%%%%%%%%%%
%%%%%%%%%%%%%%%%%%%%
\section{Cores}

A \emph{p-retraction} of $\N$ onto $\A$ is a total p-morphism $f:\N\to\A$ that is (equivalent to) the identity on $\A$.
Retractions are exactly the idempotent elements of ${\rm End}(\N)$.
In this section we characterize the minimal idempotents of ${\rm End}(\N)$, that is, the p-morphisms $u\in{\rm End}(\N)$ such that ${\rm End}(\N){*}u$ is a minimal left ideal.
These are the retractions onto substructures of $\N$ named \textit{cores.}

% \begin{fact}\label{fact_reduced_funct} 
%   Let $R$ be a p-morphism and let $f\subseteq R$ be a reduced p-morphism.
%   Then there is a reduced p-morphism $f'$ extendiung $f$ and such that 

%   \ceq{\hfill\big(\mathrel{\equiv_{\Delta,M}}\circ \mathrel{R}\circ\mathrel{\equiv_{\Delta,M}}\big)}{=}{\big(\mathrel{\equiv_{\Delta,M}}\circ \mathrel{f'}\circ\mathrel{\equiv_{\Delta,M}}\big).}
% \end{fact}

% \begin{proof}
%   Let $R$ be a p-morphism.
%   Pick a reduced p-morphism $f$ such that $R\ \subseteq\ \mathrel{\equiv_{\Delta,M}}\circ \mathrel{f}\circ\mathrel{\equiv_{\Delta,M}}$.
%   Let $\bar a$ be an enumeration of ${\rm dom} f$ and let $b\in\U^x$.
%   Let $p(\bar x)=\textrm{p-tp}(\bar a)$ and let $q(\bar x,y)=\textrm{p-tp}(\bar a,b)$.
%   By $\infty$-p-saturation and Fact~\ref{fact_p_qe}, $p(\bar x)$ is equivalent to $\exists y\,q(\bar x,y)$.
%   Then $q(f\bar a,y)$ is consistent so $f$ can be extended to a p-morphism defined on $b$.
%   The second claim follows by back-and-forth.
% \end{proof}

A set $A\subseteq\N$ is \emph{p-maximal} if it satisfies the following (obviously) equivalent conditions for every $\bar a\in A^{<\omega}$ and every p-morphism $f$ defined on $\bar a$:

\begin{itemize}
  \item[a.] $r_{\phi}(\,\bar a\,)\ \ \Leftrightarrow\ \ r_{\phi}(\,f\!\bar a\;)$ \hfill
  for every $r_{\phi}\in L^{\rm\! p}$

  \item[b.] $\textrm{p-tp}(\,\bar a\,)\ \ =\ \ \textrm{p-tp}(\,f\!\bar a\;)$.
\end{itemize}

% By Zorn's lemma, every set satisfying the above conditions is contained in a p-maximal set.

\begin{fact}\label{fact_pmaximal}
  The following are equivalent
  \begin{itemize}
    \item[1.] $A\subseteq\N$ is p-maximal
    \item[2.] $p(\bar x)=\textrm{p-tp}(\bar a)$ is a maximally consistent p-type for every $\bar a\in A^{\bar x}$ and every finite $\bar x$
    \item[3.] $p(\bar y)=\textrm{p-tp}(\bar c)$ is a maximally consistent p-type for $\bar c$ any enumeration of $A$
    \item[4.] every p-morphism defined on $A$ is a p-isomorphism between $A$ and its image.
  \end{itemize}
\end{fact}

\begin{proof}
  1$\Rightarrow$2. Negate 2. Suppose $\textrm{p-tp}(\bar a)$ is not maximal for some finite $\bar a\in A^{\bar x}$.
  Pick some $\bar a'\in A^{\bar x}$ such that $\textrm{p-tp}(\bar a)\subset \textrm{p-tp}(\bar a')$.
  Then $f=\{\langle \bar a,\bar a'\rangle\}$ is a p-morphism that contradicts the p-maximality of $A$.
  
  2$\Rightarrow$3. If $\bar c$ is an enumeration of $A$ that witnesses the failure of 3, then some finite subtuple of $\bar c$ witnesses the failure of 2.

  3$\Rightarrow$4. Negate 4. Let $f$ be a p-morphism defined on $A$ that is not a p-isomorphism.
  Let $\bar a\in A^{\bar x}$ be such that  $\textrm{p-tp}(\bar a)\smallsetminus \textrm{p-tp}(f\bar a)\neq\varnothing$.
  To contradict 3, it suffices to extend $f$ to a p-morphism defined on $A$.
  This is possible by $\infty$-p-homogeneity.

  4$\Rightarrow$1. Clear.
\end{proof}

\begin{example}
  Assume that $a\in N^x$ is such that $\phi(a\,;M^z)$ is definable over ${\rm acl^{eq}}\varnothing$ for every $\phi(x\,;z)\in L$.
  Then the orbit of $a$ under ${\rm Aut}(\U/\{M\})$ is p-maximal.
\end{example}

\begin{proof}
  First consider the case when $\phi(a\,;M^z)$ is definable over $\varnothing$.
  We prove that $\{a\}$ is p-maximal.

  Let $\phi'(z)\in L$ be such that $\phi(a\,;M^z)=\phi'(M^z)$.
  Then, if we set $\phi''(x\,;z)=\phi(x\,;z)\leftrightarrow\phi'(z)$, the formula $r_{\phi''}(a)$ holds.

  If $f$ is a p-morphism then $r_{\phi''}(fa)$ also holds.
  Then $a\equiv_{\Delta,M}f\!a$ follows.
  This shows that p-morphism are trivial, i.e.\@ equivalent to the identity, and therefore p-isomorphisms.

  Now assume that $\phi(a\,;M^z)$ is definable over ${\rm acl^{eq}}\varnothing$.?????????
\end{proof}

A set $\J\subseteq\N$ is called a \emph{core} if it is p-maximal and, among the p-maximal sets, it is maximal w.r.t.\@ inclusion.

\begin{fact}\label{fact_onto}
  Let $\J$ be a core.
  Let $\A$ be a p-maximal set.
  Then every p-morphism $f:\A\to\J$ extends to a p-isomorphism between $\A$ and $\J$. 
\end{fact}

\begin{proof}
  Reason as in the proof of Fact~\ref{fact_phogeneity} with the same notation but choose $b\in\A$.
  Let $f'$ be a p-morphism defined on $\A$ that extends $f\cup\{\langle b,c\rangle\}$.
  Then $f'[\A]$, i.e.\@ the image of $\A$ under $f'$, is a p-maximal set that contains $\J$.
  By maximality $f'[\A]$ coincides with $\J$, hence $c\in\J$.
\end{proof}

By back-and-forth we obtain the following.

\begin{corollary}[ ($\infty$-p-homogeneity of the core)]\label{fact_corephogeneity}
  If $\J$ and $\J'$ are cores, then every p-morphism $f:\J\to\J'$ extends to a p-isomorphism between $\J$ and $\J'$.
\end{corollary}

A \emph{p-retraction} of $\N$ onto $\A$ is a total p-morphism $f:\N\to\A$ that is equivalent to the identity on $\A$.

\begin{theorem}
  If $\J$ and is a core, then there is a p-retraction of $\N$ onto $\J$.
\end{theorem}

\begin{proof}
  Let $\bar c$ enumerate $\N$.
  Let $\bar a$ be a tuple such that $\textrm{p-tp}(\bar a)\supseteq \textrm{p-tp}(\bar c)$ is p-maximal.
  Let $f\ =\{\langle\bar c,\bar a\rangle\}$.
  Then $f:\N\to\A$ is a p-morphism onto the p-maximal set $\A={\rm range(\bar a)}$.
  Then $f{\restriction}\J$ is a p-isomorphism from $\J$ into $\A$.
  As $\J$ is a core, $f{\restriction}\J$ is p-isomorphism between $\J$ and $\A$.
  Then $f\circ(f{\restriction}\J)^{-1}$ is a retraction of $\N$ onto $\J$.
\end{proof}

\begin{fact}
  The following are equivalent
  \begin{itemize}
    \item[1.] $a\in\bigcap\big\{\,\J\ :\ \J \textrm{ is a core}\big\}$
    \item[2.] $\phi(a\,;M^z)$ is definable over ${\rm acl^{eq}}\,\varnothing$, for every formula $\phi(x\,;z)\in L$.
  \end{itemize} 
\end{fact}

\begin{proof}
  ???
\end{proof}

\begin{fact}
  $\J$ is compact subset of $\N$.
\end{fact}

\begin{proof}
  ???
\end{proof}

\section{Esperimenti}

Recall that $\hat m$ is an enumeration of $M$ and $|z|=|\hat m|$.

The following fact is well-known.
It entails that the relation $a$ and $b$ are at distance $\le1$ in the Lascar graph is type-definable (this is the relation defined by 1 in the fact).

\begin{fact}
  The following are equivalent
  \begin{itemize}
    \item[1.] the $a\equiv_{M'}b$ for some model $M'$
    \item[2.] $a\equiv_{m'}b$ for some $m'\equiv\hat m$.
  \end{itemize}
\end{fact}

\begin{proof}
  2$\Rightarrow$1. Clear.

  1$\Rightarrow$2. Negate 2.
  Then $\alpha(z)\vdash\phi(a\,;z)\nleftrightarrow\phi(b\,;z)$ for some $\phi(x\,;z)\in L$ and some $\alpha(z)\in {\rm tp}(\hat m)$.
  As every model witnesses $\alpha(z)$, this negates 1.
\end{proof}

Let $l_1(x,x')\subseteq L$ be the type that says that $x,x'$ are at distance $\le 1$ in the Lascar graph.

\begin{fact}  
  There is a p-type defining $l_1(x,x')$.
\end{fact}

\begin{proof}
  Consider the p-type containing the formulas $r_{\theta,\alpha}(x_0,x_1)$ for every $\theta(x_0,x_1\,;z)$ of the form $\phi(x_0\,;z)\nleftrightarrow\phi(x_1\,;z)$, where $\alpha(z)$ is the type ${\rm tp}(\hat m)$ restricted to the $z$-variables that effectively occur in $\phi(x\,;z)\in\Delta$.
\end{proof}

\begin{fact}
  If $\neg l_1(a,b)$, then $fa\not\equiv_M fb$, for every p-morphism $f$.
\end{fact}

\begin{proof}
  Assume $\neg l_1(a,b)$. Then $\alpha(z)\vdash\phi(a\,;z)\nleftrightarrow\phi(b\,;z)$ for some $\phi(x\,;z)\in L$ and some $\alpha(z)\in {\rm tp}(\hat m)$. 
  Then $r_{\theta,\alpha}(a,b)$, where $\theta(x_0,x_1\,;z)$ is the formula $\phi(x_0\,;z)\nleftrightarrow\phi(x_1\,;z)$.
  Then  $r_{\theta,\alpha}(fa,fb)$ also holds.
  As $\alpha(z)$ is satisfable in $M$, we obtain $fa\not\equiv_M fb$.
\end{proof}

\begin{corollary}
  If $f:N\to\J$ is a retraction, then $l_1(a,fa)$ for every $a,b\in\N$.
\end{corollary}




%%%%%%%%%%%%%%%%%%%%%%%%%%
%%%%%%%%%%%%%%%%%%%%%%%%%%
%%%%%%%%%%%%%%%%%%%%%%%%%%
%%%%%%%%%%%%%%%%%%%%%%%%%%
%%%%%%%%%%%%%%%%%%%%%%%%%%
\section{Notes and references}

\begin{biblist}[]\normalsize

\bib{Hr}{article}{
  label={Hr},
  author={Hrushovski, Ehud},
  title={Definability patterns and their symmetries},
  journal={arXiv:1911.01129},
}\smallskip
\bib{KNP}{article}{
  label={KNP},
  author={Krupiński, Krzysztof},
  author={Newelski, Ludomir},
  author={Simon, Pierre},
  title={Boundedness and absoluteness of some dynamical invariants in model theory},
  journal={arXiv:95100.5071},
}\smallskip
\bib{Si}{article}{
  label={Si},
  author={Simon, Pierre},
  title={Notes on Hrushovski's ``Definability patterns and
their symmetries''},
  journal={},
}\smallskip
\end{biblist}
\end{document}
